\documentclass[UTF8,12pt]{ctexart}
\usepackage[a4paper,margin=24mm]{geometry}
\usepackage{amsmath,amssymb,bm,graphicx,adjustbox}
\newcommand{\fitmath}[1]{\begin{adjustbox}{max width=\linewidth}$\displaystyle #1$\end{adjustbox}}
\setlength{\parindent}{0pt}
\setlength{\parskip}{.7em}
\sloppy
\allowdisplaybreaks
\begin{document}
\section*{2007 年考研数学三 · 参考答案}
\subsection*{2007年真题参考答案}

\subsection*{一、选择题}

（1）B。（2）D。（3）C。（4）B。（5）D。（6）D。（7）A。（8）B。（9）C。（10）A。


\subsection*{二、填空题}

（11）\(\displaystyle 0\)。（12）\(\displaystyle \dfrac{\displaystyle (-1)^n2^nn!}{\displaystyle 3^{n+1}}\)。（13）\(\displaystyle -\dfrac{\displaystyle 2y}{\displaystyle x}f_1^{\prime}+\dfrac{\displaystyle 2x}{\displaystyle y}f_2^{\prime}\)。（14）\(\displaystyle \dfrac{\displaystyle x}{\displaystyle \sqrt{\ln x+1}}\)。（15）\(\displaystyle 1\)。（16）\(\displaystyle \dfrac{\displaystyle 3}{\displaystyle 4}\)。


\subsection*{三、解答题}

（17）曲线 \(\displaystyle y=y(x)\) 在点 \(\displaystyle (1,\allowbreak 1)\) 附近是凸的。（18）\(\displaystyle \dfrac{\displaystyle 1}{\displaystyle 3}+4\sqrt2\ln(\sqrt2+1)\)。（19）证明略。


（20）\(\displaystyle \dfrac{\displaystyle 1}{\displaystyle x^2-3x-4}=-\dfrac{\displaystyle 1}{\displaystyle 5}\displaystyle \sum_{n=0}^{\infty}\left[\dfrac{\displaystyle 1}{\displaystyle 3^{n+1}}+\dfrac{\displaystyle (-1)^n}{\displaystyle 2^{n+1}}\right](x-1)^n\)，\(\displaystyle x\in(-1,\allowbreak 3)\)。


（21）当 \(\displaystyle a=1\) 时，\(\displaystyle x=k(1,\allowbreak 0,\allowbreak -1)^{\mathrm T}\)，其中 \(\displaystyle k\) 为任意常数，是方程组\textcircled{1}和\textcircled{2}的公共解；当 \(\displaystyle a=2\) 时，\(\displaystyle x=(0,\allowbreak 1,\allowbreak -1)^{\mathrm T}\) 是方程组\textcircled{1}和\textcircled{2}的唯一公共解。


（22）（I）矩阵 \(\displaystyle B\) 的属于特征值 \(\displaystyle -2\) 的特征向量为 \(\displaystyle k_1(1,\allowbreak -1,\allowbreak 1)^{\mathrm T}\)，其中 \(\displaystyle k_1\) 为任意非零常数；矩阵 \(\displaystyle B\) 的属于特征值1的特征向量为 \(\displaystyle k_2(1,\allowbreak 1,\allowbreak 0)^{\mathrm T}+k_3(-1,\allowbreak 0,\allowbreak 1)^{\mathrm T}\)，其中 \(\displaystyle k_2,\allowbreak k_3\) 为不全为零的常数。（II）\(\displaystyle B=\begin{pmatrix}0&1&-1\\1&0&1\\-1&1&0\end{pmatrix}\)。


（23）（I）\(\displaystyle \dfrac{\displaystyle 7}{\displaystyle 24}\)。（II）\(\displaystyle f_Z(z)=\begin{cases}2z-z^2,\allowbreak &0<z\le1,\allowbreak \\z^2-4z+4,\allowbreak &1<z\le2,\allowbreak \\0,\allowbreak &\text{其他}.\end{cases}\)


（24）（I）\(\displaystyle \hat\theta=2\bar X-\dfrac{\displaystyle 1}{\displaystyle 2}\)。（II）\(\displaystyle E(4\bar X^2)\ne\theta^2\)。因此，\(\displaystyle 4\bar X^2\) 不是 \(\displaystyle \theta^2\) 的无偏估计量。

\end{document}
